一个\emph{函数}~$f\colon A \to B$ 将
\emph{定义域}~$A$ 中的每一个元素映射到
\emph{陪域}~$B$ 中唯一的一个元素。
若 $x \in A$，则 $f$ 将 $x$ 映射到的那个 $y$ 称为
$f$ 在\emph{自变元}~$x$ 处的\emph{值}~$f(x)$。
如果 $A$ 是一个序对的集合，我们可以认为函数~$f$ 接受两个自变元。
$f$ 的\emph{值域}~$\ran{f}$ 是~$B$ 的一个子集，由~$f$ 的所有值组成。

如果 $\ran{f} = B$，那么 $f$ 称为\emph{满射的}。
值~$f(x)$ 是唯一的，即 $f$ 将 $x$ 只映射到一个 $f(x)$，绝不会更多。
如果 $f(x)$ 在另一种意义下也是唯一的，即没有两个不同的自变元被映射到相同的值，那么 $f$ 称为\emph{单射的}。
既是单射又是满射的函数称为\emph{双射的}。

双射函数有唯一的\emph{逆函数}~$f^{-1}$。
函数也可以组合在一起：函数 $(g \circ f)$ 是 $f$ 与 $g$ 的\emph{复合}。
单射函数的复合是单射的，满射函数的复合是满射的，并且 $(f^{-1} \circ f)$ 是恒等函数。

如果我们放宽 $f$ 必须对每个 $x \in A$ 都有值的要求，就得到了\emph{偏函数}的概念。
如果 $f\colon A \pto B$ 是偏函数，那么当 $f$ 对自变元~$x$ 有值时，我们说 $f(x)$ 是\emph{有定义的}，记作 $f(x) \fdefined$；否则我们说 $f(x)$ 是\emph{无定义的}，记作 $f(x) \fundefined$。
任何（偏）函数~$f$ 都关联着一个\emph{图像}~$R_f$，即当且仅当 $f(x) = y$ 时成立的关系。